% !TEX root = ../main.tex

\chapter*{Abstrak}
\addcontentsline{toc}{chapter}{Abstrak}

Relativitas Umum memprediksi bahwa massa berotasi menginduksi penyeretan kerangka acuan inersia lokal, yang dikenal sebagai efek Lense-Thirring atau \textit{frame-dragging}.
Efek ini telah dikonfirmasi secara eksperimental oleh misi \textit{Gravity Probe B} (GP-B) dengan hasil pengukuran sebesar $37{,}2 \pm 7{,}2$ milliarcsecond per tahun (mas/yr) untuk presesi \textit{frame-dragging} dan $6602 \pm 18$ mas/yr untuk presesi geodesik.
Penelitian ini menyajikan kajian analitik terhadap persamaan geodesik dan transportasi paralel tensor spin pada metrik Kerr, serta penyelesaian numeriknya untuk konfigurasi orbit Bumi.
Metrik Kerr dalam koordinat Boyer-Lindquist diturunkan secara eksplisit dalam bentuk matriks $4 \times 4$, dan simbol-simbol Christoffel yang relevan dihitung untuk memperoleh persamaan gerak.
Lagrangian metrik Kerr digunakan untuk menurunkan persamaan geodesik melalui metode Euler-Lagrange, yang kemudian direduksi menjadi sistem persamaan diferensial biasa orde satu.
Aproksimasi medan lemah diterapkan untuk memperoleh ekspresi analitik laju presesi Lense-Thirring $\Omega_{\text{LT}} \approx 2GJ/(c^2 r^3)$.
Sistem persamaan diferensial yang dihasilkan diselesaikan secara numerik menggunakan metode Runge-Kutta orde empat yang diimplementasikan dalam Python.
Hasil simulasi menunjukkan laju presesi geodesik sebesar $\sim\!6600$ mas/yr dan \textit{frame-dragging} sebesar $\sim\!37$ mas/yr, konsisten dengan data eksperimen GP-B dalam batas galat pengukuran.
Analisis kestabilan numerik dan sensitivitas terhadap parameter rotasi Bumi turut disajikan untuk memvalidasi keandalan metode komputasi yang digunakan.

\vspace{1em}
\noindent\textbf{Kata Kunci:} Relativitas Umum, Metrik Kerr, Efek Lense-Thirring, \textit{Frame-Dragging}, Presesi Geodesik, Gravity Probe B, Penyelesaian Numerik.
